\section{The canonical weak splitting map}\label{sec:introduction}

Fix a prime $p$ and work in ordinary $p$-local spectra, with sphere
$S=S^0_{(p)}$. Let $K(a)$ denote Morava $K$-theory of height $a$, put
$K(0)=H\Q$, and write
\[
 L_r=L_{K(0)\vee\cdots\vee K(r)},\qquad
 X=S_p^\wedge=\operatorname{holim}_{j\geq1}S/p^j.
\]
Thus $L_r$ is Johnson--Wilson $E(r)$-localization. For a spectrum $Z$,
let $\eta_Z^{(n)}:Z\to L_{K(n)}Z$ be the localization unit. We study
\begin{equation}\label{eq:weak-unit}
 i_{n,X}=L_{n-1}(\eta_X^{(n)}):
 L_{n-1}X\longrightarrow L_{n-1}L_{K(n)}X.
\end{equation}
The weak chromatic splitting conjecture predicts a homotopy retraction
of this map whenever $X$ is the derived $p$-completion of a finite
spectrum. A retraction is a map of spectra $r$ with
$r i_{n,X}\simeq\mathrm{id}$; it need not be natural in the finite input.
We disprove this universal assertion with finite input the sphere.

For an odd prime, let
\[
 \beta_1\in\pi_{d_p}S,\qquad
 d_p=2p(p-1)-2,\qquad k_p=(p-1)^2,
 \qquad z_p=\beta_1^{k_p},
\]
where $\beta_1$ is the classical first beta element.

\begin{theorem}[A named kernel at height $p$]
\label{thm:height-p-kernel}
For every prime $p\geq5$, the image of $z_p$ in
$\pi_{d_pk_p}L_{p-1}X$ is nonzero and is killed by
$\pi_{d_pk_p}(i_{p,X})$. Consequently $i_{p,X}$ has no homotopy
retraction.
\end{theorem}

\begin{theorem}[The same kernel at height $p+1$]
\label{thm:height-plus-one-kernel}
For every prime $p\geq7$, the image of the same $z_p$ in
$\pi_{d_pk_p}L_pX$ is nonzero and is killed by
$\pi_{d_pk_p}(i_{p+1,X})$. Consequently $i_{p+1,X}$ has no homotopy
retraction.
\end{theorem}

\begin{corollary}[The prime-five class]\label{thm:five}
At $p=5$, the image of $\beta_1^{16}$ in $\pi_{608}L_4X$ is nonzero
and belongs to the kernel of $\pi_{608}(i_{5,X})$.
\end{corollary}

The proof detects $z_p$ under a height-$(p-1)$ unit and kills it under
both upper units. The lower detection follows from the Hopkins--Miller
calculation recorded by Heard
\cite[proof of Theorem~5.6 and Lemma~5.8]{heard2015}: the actual unit
detects the named $\beta_1$, and its $k_p$th power remains nonzero in
the ordinary homotopy fixed point spectral sequence. At the upper
height $a$, we construct a separated multiplicative filtration with
the unit image of $\beta_1$ in filtration at least two and
$F^{a^2+2}=0$. The required
inequalities are
\[
 2k_p-(p^2+1)=p^2-4p+1>0\quad(p\geq5),
\]
\[
 2k_p-((p+1)^2+1)=p(p-6)>0\quad(p\geq7).
\]
Completion and naturality then place the nonzero lower image in the
kernel of the exact map \eqref{eq:weak-unit}. For example, at $p=7$
the class $\beta_1^{36}$ has degree $82\cdot36=2952$ and gives
kernels at both heights seven and eight.

\paragraph{Notation.}\label{par:common-notation}
Let $E_a$ be Morava $E$-theory of the height-$a$ Honda formal group
over $\F_{p^a}$, and put $B_a=L_{K(a)}S$. Unless a localized tensor
is explicitly displayed, smash products and function spectra are
ordinary. All homotopy groups are those of the underlying spectra.

\subsection{An independent descent-exponent obstruction}

The height-$p$ nonretraction also follows from incompatible finite
descent lengths. Put $t=p-1$, and let the finite pure-stabilizer subgroup
$H=C_p\rtimes C_{(p-1)^2}$ act on $E_t$. In
$\mathcal D_H=\Fun(BH,\Sp)$, with pointwise ordinary smash, put
$A_H=F(H_+,S)$ with its translation action and
$J_H=\fib(S\to A_H)$. Define $\expn_H(V)$ as the least $b\geq1$
such that $J_H^{\wedge b}\wedge V\to V$ is null, and as infinity
if no such $b$ exists. Set
\[
 T(V)=L_{K(t)}(E_t\wedge V),
\]
with $H$ acting on the first factor.

\begin{theorem}[Incompatible descent exponents]\label{thm:exponents}
For every prime $p\geq5$, with $m=p^2+2$,
\[
 \expn_H(E_{p-1})>m,
 \qquad \expn_H\bigl(T(B_p)\bigr)\leq m.
\]
A homotopy retraction of $i_{p,X}$ would make $E_{p-1}$ an
$H$-equivariant retract of $T(B_p)$. Hence $i_{p,X}$ has no homotopy
retraction, independently of Theorem~\ref{thm:height-p-kernel}.
\end{theorem}

For the lower bound, a class survives to one specified finite page of
the ordinary homotopy fixed point spectral sequence. No permanent
cycle or sphere representative is needed. For the upper bound,
Devinatz--Hopkins' nilpotence theorem supplies a finite exponent before
any homotopy test is chosen \cite[Appendix~A, Proposition~A.3]{dh2004}.
The numerical cochain bound then makes the relevant map zero after
every finite ordinary test. A Brown--Comenetz argument proves that
the map itself is null, and a coherent trace construction transports
the resulting finite filtration to height $p-1$. The finite-stage
formalism used here belongs to the nilpotence and descent theory of
Mathew and Mathew--Naumann--Noel
\cite{mathew2016,mnn2017,mnn2019}.

\subsection{Relation to earlier splitting results}

Hopkins' chromatic splitting conjecture, recorded by Hovey
\cite[Conjecture~4.2, especially part~(v)]{hovey1995}, predicts a
specific decomposition of chromatic overlap as well as the weak unit
retraction. These predictions differ already at $n=p=2$.
Beaudry disproves the proposed strong summand formula
\cite[Theorem~1.0.4]{beaudry2017}; Beaudry--Goerss--Henn give a revised
decomposition with Moore terms that retains the canonical unit as a
retract \cite[Theorem~1.1.6]{bgh2022}. Their
Conjectures~1.1.10--1.1.11 distinguish the strong and
completed-finite-input weak formulations, and their introduction
records the weak form at all primes through height two. The theorem
numbers refer to the versioned texts identified in the bibliography.
Outside finite inputs, Minami reports Devinatz's counterexample for
a $BP$ analogue allowing completed $BP$
\cite[Introduction and Remark~1.1]{minami2003}. The present
counterexamples use the completion of the finite sphere.

Long beta powers in the sphere were known earlier. Lee--Ravenel prove
$\beta_1^{p^2-p-1}\ne0$ for $p\geq7$ and report
$\beta_1^{17}\ne0$, $\beta_1^{18}=0$ at five
\cite[opening theorem and introduction]{lee-ravenel1994}.
Ravenel's first-beta detection calculation relates the cobar class to
stabilizer cohomology \cite[Theorem~4 and its proof in Section~2]{ravenel1978}.
The additional requirement here is detection after a specified lower
localization and vanishing of that same named sphere product under
the upper unit.

\paragraph{Organization.}
Sections~\ref{sec:completion}--\ref{sec:kernels} prove the named
kernel results. Their upper bound uses dimension at height $p+1$
and a semilinear trace-one contraction at height $p$, retaining the
possible inverse-limit degree in both cases.
Sections~\ref{sec:towers}--\ref{sec:exponent-transport} give the
independent exponent proof. Appendix~\ref{sec:ordinary-products}
proves the ordinary filtered-product lemma used by the lower detector
and by the finite-page argument.
